\subsection{Cancelación uniforme entre inserción y marginal corregida}\label{sec:B4T09}
Mantenemos las hipótesis, las normalizaciones y la rama pequeña de las
Secciones~\ref{sec:B4T01}--\ref{sec:B4T08}.
La corrección parental es $\delta_c=(R\nabla\ell)_c/p_c$.
No se identifica la marginal corregida con la marginal verdadera del conteo.

\paragraph{Derivada del prefactor en las variables de diagonal uno.}
La identidad exacta de la Sección~\ref{sec:B4T03} y la identidad de la resolvente para $S=S_0+E$ dan
\[
 G_{ab}-\Gamma_{ab}=-3m_{2,a}\mathbf1_{a=b}+O(k^{-2}).
\]
En efecto, $E_{aa}=3m_{2,a}+O(k^{-2})$, $E_{ab}=O(k^{-2})$
fuera de la diagonal, $\|E\|_\infty=O(k^{-1})$,
$\Gamma=I+O_{\max}(k^{-1})$ y $\|\Gamma\|_\infty=O(1)$.
Al expandir la fórmula exacta de $g=\nabla\ell$, resulta
\[
 g_i+1=-6p_i-f_i+\xi_{r_i}+\zeta_{c_i}+O(k^{-2}),
\]
donde se puede tomar
\[
 \xi_r=-\tfrac12(\Gamma_{rr}-1)+\tfrac32m_{2,r},\qquad
 \zeta_c=-\tfrac12(\Gamma_{cc}-1)+\tfrac32m_{2,c}.
\]
Las funciones RC se anulan al aplicar la respuesta tangente.
Además, con $R_0=P_p-P_pA\Gamma A^TP_p$,
$\|P_p^{-1}(R-R_0)\|_\infty=O(k^{-1})$ por la identidad de la resolvente.
Como $-6p-f=O(k^{-1})$, se puede sustituir $R$ por $R_0$ con precisión
relativa $O(k^{-2})$. Por tanto
\[
 \delta_c=-6p_c-f_c+
 \sum_i p_i(a_i^T\Gamma a_c)(6p_i+f_i)+O(k^{-2}).
\]
En $F\setminus c$, $a_i^T\Gamma a_c=1+O(k^{-1})$.
En $E'$, la Sección~\ref{sec:B4T06} da $p_i a_i^T\Gamma a_c=-y_i+O(p_i/k^2)$.
El término excepcional de $c$ sólo produce un cambio $O(k^{-2})$, así que
\[
 \boxed{\delta_c=-6p_c-f_c+
 \sum_Fp_i(6p_i+f_i)-\sum_{E'}y_i(6p_i+f_i)+O(k^{-2}).}
\]
Todas las sumas de errores respetan $\sum_Ep=k$ y $\sum_Fp=O(1)$.

\paragraph{Comparación con el hijo.}
Sea $\epsilon_c=L(B+c)-L(B)-\log p_c$.
Combinar la energía de la Sección~\ref{sec:B4T06}, la sustitución de la Sección~\ref{sec:B4T07} y
el desarrollo del determinante de la Sección~\ref{sec:B4T08} da
\[
 \begin{aligned}
 \epsilon_c={}&-6p_c-f_c+6\sum_Fp_i^2-6\sum_{E'}p_iy_i
 -\sum_{E'}y_i f_i\\
 &+\sum_Fp_i f_i^I-\sum_{j\in J}(K_0)_{jj}-1/k+O(k^{-2}).
 \end{aligned}
\]
La tabla siguiente indica de dónde procede cada término; en la segunda fila se usa
$a_c^T\Gamma a_c-4=\sum_{j\in J}(K_0)_{jj}+2f_c$.
\begin{center}
\renewcommand{\arraystretch}{1.3}
\begin{tabular}{@{}ll@{}}
\toprule
Origen & Aportación a $\epsilon_c$, salvo $O(k^{-2})$\\
\midrule
Frontera de la acción (Sección~\ref{sec:B4T04}) & $\tfrac32\sum_Fp_i^2$\\
Energía $-\tfrac12\mathcal Q(y_0)$ (Sección~\ref{sec:B4T06}) &
 $-\tfrac12\sum_{j\in J}(K_0)_{jj}-f_c-6p_c$\\
Factor local (Sección~\ref{sec:B4T07}) & $3\sum_Fp_i^2-6\sum_{E'}p_iy_i$\\
Determinante (Sección~\ref{sec:B4T08}) &
 $-\tfrac12\sum_{j\in J}(K_0)_{jj}-\tfrac1k+\sum_Fp_if_i^I-\sum_{E'}y_if_i+\tfrac32\sum_Fp_i^2$\\
\bottomrule
\end{tabular}
\end{center}
Los términos $\tfrac12\log(2\pi)$ coinciden porque ambos soportes
son RC conexos; el covolumen ya está incorporado mediante cofactores.
Al restar la fórmula de $\delta_c$, queda
\[
 \epsilon_c-\delta_c=-\sum_{j\in J}(K_0)_{jj}-1/k
                    -\sum_Fp_i(f_i-f_i^I)+O(k^{-2})=O(k^{-2})
\]
por la identidad final de la Sección~\ref{sec:B4T08}.
Como $\delta_c=O(k^{-1})$,
\[
 \boxed{r_c=L(B+c)-L(B)-\log[p_c(1+\delta_c)]=O(k^{-2}).}
\]
Las constantes dependen sólo de $\delta_0,C,\eta,\gamma$, y existe
$k_0(\delta_0,C,\eta,\gamma)$ tal que el resultado vale para todo
$k\ge k_0$ y todo hijo legal del padre, sin selección de hijos.

\paragraph{Alcance.}
El punto estacionario hijo usado es el de la rama $p<1/4$, cuya existencia y unicidad
prueba la Sección~\ref{sec:B4T02}. No se afirma que las hipótesis cuantitativas valgan en todos
los residuales del problema de conteo.
